\section{Certification of the mixed-cover construction}\label{app:cover}
This appendix specifies the hard case in Lemma~\ref{lem:mixed}. The
certificate concludes the existence of actual packings; it is not merely
a numerical upper bound for a hypothetical sample.

\subsection{Coordinates and fixed-size samples}
For $m>0$ put
\[
 x=\frac{\sqrt n}{B_*m^{1/3}},\quad s'=\frac{n^2}{2m}\ge1,\quad
 W(x)=B_*^3x^2,\quad v(x)=W(x)(1+3x).
\]
Then
\begin{equation}\label{eq:cover-coordinates}
 n=\frac{4B_*^3(s')^2}{W(x)^3},\qquad
 m=\frac{n^2}{2s'},\qquad D(n,m)/m=v(x).
\end{equation}
The irreducible range is $0<x\le17/250$, $s'\ge1$.
Set
\[
 p_0(x)=v(x)\left(\frac{99}{200}+\frac{11}{100}v(x)
                                      +\frac7{100}v(x)^2\right),\qquad
 N=\floor{p_0(x)n},\quad p=N/n.
\]
The sample is uniform among the $N$-element subsets of $V(G)$.
In particular, the proof does not substitute independent vertex sampling
for this fixed-cardinality law. For $P=p/W(x)$ and
$P_0(x)=p_0(x)/W(x)$ the exact floor bounds are
\[
 P_0(x)-\frac{W(x)^2}{4B_*^3(s')^2}\le P\le P_0(x).
\]
Monotone endpoint bounds turn these into rational bounds on every parameter
rectangle.

Let
\[
 H(x)=103/250+(351/200)x,\qquad D_0(x)=1-(33/125)W(x),
 \]\[
 h_G=B_*m^{1/3}H(x)/D_0(x),\qquad z=ph_G.
\]
Here $H(x)$ is a local coordinate function, unrelated to the loss $H(w)$
in Section~\ref{sec:potential}. The graph is assumed to meet the sparse-cut
irreducibility and minimum-degree inequalities left by the induction in
Lemma~\ref{lem:mixed}.

A good sample $U$ has size exactly $N$, certified lower neighbour degrees
into $U$, cut expansion inside $U$, an anchoring inequality, at least
$(97/100)p^2m$ induced edges, and at most $D(n,m)$ edges incident with $U$.
The anchoring inequality, for $X\subseteq V\setminus U$ and
$Y\subseteq U$ with $2\abs Y\le\abs U$, is
\[
 L\abs X\le e_G(X,U\setminus Y)+e_G(Y,U\setminus Y).
\]
Some rectangles also require a local cut bound up to sample proportion
$1/125$. The certified selection theorem bounds the union of six bad events:
low degree, a failed global cut, failed anchoring, too few sample edges,
too many incident edges, and the optional local-cut failure. Their rational
upper bounds have sum strictly less than $1$. Thus at least one actual
$N$-element sample is good. The fixed-size moment, degree, marked-cut and
anchor bounds and their parameter matches are proved before this union
bound is used.

\subsection{From the sample to cycles}
A capacity certificate splits the induced sample edges into banks, constructs
a weak clique immersion in the appropriate bank, and uses the remaining
expansion to route the requests from the path-and-cycle decomposition of Lov\'asz~\cite{Lovasz1968} applied to the
complement. Weak immersion is sufficient: routed paths are edge-disjoint;
branch vertices are not silently assumed to be avoided by all interiors.
The indexed routing and cycle-closure lemmas account for endpoint occurrences
and integer terminal capacities. The global-cut and additional-local-cut
branches have separate capacity inequalities.

The output is an actual packing $\mathcal C$ in $G$ with
\begin{equation}\label{eq:sample-output}
 p(\mathcal C)\le\frac{n-\abs U}{2},\qquad
 q(\mathcal C)\le D(n,m)+n-\abs U.
\end{equation}
The sample-size lower bound, including its floor, pays the fine cycle budget
$n\psi(m/(2B_*^3n))$. The second inequality in~\eqref{eq:sample-output}
pays both allowed residual-edge budgets, and the first is also at most
$n/2$. Thus this hard case proves both conclusions of
Lemma~\ref{lem:mixed}. The full induction may still choose different
packings for its two conclusions.

\subsection{Continuous coverage and its soundness interface}
A rational rectangle is a tuple $(x_-,x_+,s_-,s_+)$, where $s_+$ may be
absent, denoting an unbounded ray. Its \emph{soundness} statement quantifies
over every actual graph whose coordinates lie in it and which satisfies
the irreducibility and degree conditions, and concludes
\eqref{eq:fine-cover} and~\eqref{eq:coarse-cover}.
A routed case consists of a sample certificate and a capacity certificate.
Their common parameters---degree, expansion, anchoring and any optional
local cut constant---must agree exactly. The theorem
\mcode{rat_routed_case_sound} combines their graph-theoretic conclusions
only after these equalities and domain inclusions are checked.

Rectangles are joined by exact rational interval chains. Horizontal joining
checks that the $x$ intervals cover and that each piece extends over the
required $s'$ interval. Vertical joining checks the converse. Both operations
include closed endpoints, and an unbounded target ray must be covered by
a chain that actually reaches an unbounded piece. These are proved
soundness-preserving operations, rather than a sampled-grid argument.
The final assembly is
\[
 [0,3/1000]\times[1,\infty)\quad\cup\quad
 [3/1000,17/250]\times[1,\infty),
\]
with each rectangle further split into its certified scale pieces.
Its final statement is
\begin{center}\mcode{Cycle12.CoverAssembly.Full.sound}.\end{center} Actual graphs have $x>0$;
including $x=0$ in the interval certificate causes no division by zero in
the graph application.

This closes the irreducible case and hence the strong induction. The
certificate hierarchy contains the actual probability bounds and routing
capacities, not just the final interval endpoints. Its rational decisions
are checked in Lean, and its probabilistic and graph-theoretic soundness
lemmas lie in the same axiom-audited dependency closure as the main theorem.
